\documentclass[12pt,a4paper]{article}
\usepackage{\string"../../../../assets/preamble"}

\begin{document}

% Section
\section*{\huge{Merge Sort}}
\vspace{.5cm}

% section
\section{Description}
Recursive algorithm based on the divide and conquer strategy. It is based on two components: the  function that split the given list recursively until it receives a single item list, and a merge function that is used on the recursive algorithm to merge two already sorted lists into a combined sorted one.

% section
\section{Time Complexity}
\vspace{.5cm}
\begin{center}

  \begin{tikzpicture}
    \begin{axis}[
        axis_style,
        title={},
      ]
      \addplot[
        main_plot_style,
        domain=0:28, % Range of the logarithmic function,
      ] {6*x*log2(x)+6*x} node[node_style]{\(6n*log_2n+6n\)};
      \addplot[
        background_plot_style,
        domain=0:25,
      ] {x^2} node[node_style]{\(x^2\)};
      \addplot[
        background_plot_style,
        domain=0:10,
      ] {2^x} node[node_style]{\(2^n\)};
      \addplot[
        background_plot_style,
        domain=0:100, % Range of the logarithmic function
      ] {x*log2(x)} node[node_style]{\(n^2\)};
      \addplot[
        background_plot_style,
        domain=0:100, % Range of the logarithmic function
      ] {log2(x)} node[node_style,above]{\(log_n\)};
    \end{axis}
  \end{tikzpicture}
\end{center}

\section{Pseudocode}
\begin{algorithmic}
  \State \textbf{Input:} \(A = [A_1,A_2,...,A_n]\)
  \State \textbf{Output:} \(B = [B_1,B_2,...,B_n]\)
  \State \noindent\rule{\textwidth}{0.4pt}
  \smallskip
  \Function{MergeSort}{$A$} \Comment{Algorithm itself}
  \State $n \gets length(A)$;
  \If{$n\leq 1$}:
  \State \Return $A$; \Comment{Length 1, already sorted, return}
  \Else
  \State $l \gets [A_1,...,A_{n/2}]$; \Comment{Split left subarray}
  \State $r \gets [A_{(n/2)+1},...,A_n]$; \Comment{Split right subarray}
  \medskip
  \State $L \gets \Call{MergeSort}{l}$; \Comment{Sort recursively}
  \State $R \gets \Call{MergeSort}{r}$; \Comment{Sort recursively}
  \State $B \gets \Call{Merge}{L,R}$; \Comment{Merge both subarrays}
  \medskip
  \State \Return $B$;
  \EndIf
  \EndFunction
  \bigskip
  \setcounter{ALG@line}{0}
  \Function{Merge}{$L, R$}
  \State $n \gets length(L) + length(R)$; \Comment{Full length of resulting array}
  \State $B \gets [\enspace]$; \Comment{Initialize empty array}
  \State $i \gets 0$;
  \State $j \gets 0$;
  \medskip
  \For{$k \gets 0$ to $n$} \Comment{Iterate full length of both subarrays}
  \If{$i \leq length(L)$ \textbf{and} $j \leq length(R)$} \Comment{Both within range}
  \If{$L[i]\leq R[j]$}: \Comment{Left lower, set it}
  \State $B[k] \gets L[i]$;
  \State $i \gets i + 1$;
  \ElsIf{$L[i] > R[j]$}: \Comment{Right lower, set it}
  \State $B[k] \gets R[j]$;
  \State $j \gets j + 1$;
  \EndIf
  \ElsIf{$i + 1 \leq length(L)$}: \Comment{Only left within range, set it}
  \State $B[k] \gets L[i]$;
  \State $i \gets i + 1$;
  \ElsIf{$j + 1 \leq length(R)$}: \Comment{Only right within range, set it}
  \State $B[k] \gets R[j]$;
  \State $j \gets j + 1$;
  \EndIf
  \EndFor
  \State \Return $B$;
  \EndFunction
\end{algorithmic}

\vfill
\begin{flushright}
  \rule{4cm}{0.4pt}\\
  \textit{\small{Antonio Díaz Correa}}\\
  \textit{\small{\today}}
\end{flushright}


\end{document}
